% Fichier GÉNÉRÉ (v0.21.0) — les comptes EXACTS des journaux
% récupérés, la source unique de vérité partagée par les cinq pièces.
% (régénéré par scripts/vague10/regenerer_macros_v210.py : les comptes
% portés aux journaux APRÈS la vague 10 + les canons de la vague 10)
\newcommand{\nbLignesJournal}{13,648}
\newcommand{\nbLignesWorklog}{13,988}
\newcommand{\nbLignesDecisions}{4,814}
\newcommand{\plageDecisions}{D1--D234}
% LES DOI (le concept résout sur la dernière version)
\newcommand{\doiConcept}{10.5281/zenodo.24367343}
\newcommand{\doiVersionCourante}{10.5281/zenodo.24382460}
% LES CANONS M4 (le record daté de la vague 9 : n=5, la liste fixe, les
% médianes RECALCULÉES des JSON bruts — le ratio de régime courant
% RECALCULÉ des médianes non arrondies : 6.912/2.056 =
% 3.362 → 3.4 à une décimale, PARTOUT où le
% courant est cité)
\newcommand{\mquatreAvantMediane}{39.3}
\newcommand{\mquatreChaudMediane}{6.9}
\newcommand{\mquatreBrasC}{2.1}
\newcommand{\mquatreRatioAvant}{19.1$\times$}
\newcommand{\mquatreRatioChaud}{3.4$\times$}
% LE COMPTE D'ARBRES SIGNÉS (le sceau I — quatorze, D177)
\newcommand{\nbArbresSignes}{14}
% LES CANONS DE LA VAGUE 9 (reconduits — le record daté)
\newcommand{\sureteDecharges}{2006}
\newcommand{\sureteDechargesTotal}{2098}
\newcommand{\sureteNOMMES}{92}
\newcommand{\mutationBrutNum}{811}
\newcommand{\mutationDenomBrut}{832}
\newcommand{\mutationAjusteNum}{811}
\newcommand{\mutationDenomAjuste}{817}
\newcommand{\cfTroisObligations}{14/14}
% LES CANONS DE LA VAGUE 10 (les pièces des ateliers, les médianes des
% JSON bruts)
\newcommand{\cgEffetNotationMs}{37073}
\newcommand{\cgEffetNotationRatio}{18.1$\times$}
\newcommand{\cgEffetGouvernanceMs}{215}
\newcommand{\cgEffetGouvernanceRatio}{1.11$\times$}
\newcommand{\cgMquatreT}{39.24}
\newcommand{\cgMquatreC}{1.95}
\newcommand{\cgMquatreCG}{2.16}
\newcommand{\fenetreEditerTotal}{1146}
\newcommand{\fenetreEditerReponse}{40}
\newcommand{\fenetreHdeuxTotal}{2452}
\newcommand{\fenetreReduction}{53.3}
\newcommand{\sseDeuxVdeuxInstructions}{4087}
\newcommand{\sseDeuxTempsReel}{833}
\newcommand{\gccODeuxTempsReel}{16040}
\newcommand{\idiomeMemchrTempsReel}{561}
\newcommand{\gardesNombre}{58}
\newcommand{\gardesLea}{34}
\newcommand{\niBlocsCouverts}{13}
\newcommand{\niFamilles}{7}
